\input{solution-sheet.tex}
\title{Unit 09 solutions}
\author{Verified by \texttt{verify/check\_math.py}}
\begin{document}
\maketitle
\begin{description}
\item[\texttt{U09-D01}] \(\pi/2\). Power-reduce.
\item[\texttt{U09-D02}] \(\pi/4\). Power-reduce.
\item[\texttt{U09-D03}] \(1/4\). $u=\cos x$.
\item[\texttt{U09-D04}] \(1-\pi/4\). $\tan^2=\sec^2-1$.
\item[\texttt{U09-D05}] \(\ln(1+\sqrt{2})\). Trig/hyp sub.
\item[\texttt{U09-D06}] \(\frac12\ln(1+\sqrt{2})+\frac{\sqrt{2}}{2}\). Standard formula.
\item[\texttt{U09-D07}] \(2/3\). $\sin^3=(1-\cos^2)\sin$.
\item[\texttt{U09-D08}] \(\pi/3\). $\arctan\sqrt3=\pi/3$.
\end{description}
\subsection*{Compact answer\_sympy}
\begin{verbatim}
U09-D01 pi/2
U09-D02 pi/4
U09-D03 1/4
U09-D04 1-pi/4
U09-D05 log(1+sqrt(2))
U09-D06 log(1+sqrt(2))/2 + sqrt(2)/2
U09-D07 2/3
U09-D08 pi/3
\end{verbatim}
\subsection*{Not a drill}
The early half-angle pass $\int_0^{\pi/3}(1+\cos x)^{-1}\,dx=\sqrt{3}/3$ is taught in \texttt{teaching.tex}. Its antiderivative is $\tan(x/2)$. It is not \texttt{U09-D01} through \texttt{U09-D08}.
\end{document}
